\chapter{Enhanced Liver Fibrosis (ELF) Test}

\section*{What Is This Test?}
The Enhanced Liver Fibrosis test is a blood-based score that combines markers of extracellular-matrix remodeling—hyaluronic acid, PIIINP, and tissue inhibitor of metalloproteinase-1—to estimate the likelihood of clinically significant liver fibrosis.

\section*{Alternative Names}
ELF Test; Enhanced Liver Fibrosis Score; Serum Liver Fibrosis Panel.

\section*{Why This Test Is Ordered}
ELF may help risk-stratify people with chronic liver disease, metabolic dysfunction-associated steatotic liver disease, viral hepatitis, or other conditions in which advanced fibrosis is a concern. It can help identify patients who need specialist assessment, but it does not replace clinical evaluation or every other fibrosis assessment.

\section*{How This Test Is Performed}
A blood sample is collected and the three fibrosis-related biomarkers are measured. The laboratory combines them into a validated ELF score, which is interpreted with liver enzymes, platelet count, disease cause, and other noninvasive assessments.

\section*{Preparation, Risks, and Limitations}
Follow local instructions about fasting. Venipuncture is the usual risk. Acute inflammation, cholestasis, age, kidney disease, and non-liver conditions can affect components of the score. Cutoffs and approved uses vary by country and laboratory.

\section*{Understanding Results}
A higher score indicates a greater likelihood of advanced fibrosis and future liver-related complications. A low score can be reassuring in the appropriate population, but neither result alone establishes cirrhosis or excludes all clinically important disease.

\section*{References}
European Association for the Study of the Liver guidance on noninvasive tests for liver disease.

\clearpage
\section*{Understanding Results and Follow-up}
The laboratory report should identify the specimen, method, units, reference interval or decision limit, and whether the result is positive, negative, abnormal, or inconclusive. Reference intervals vary with age, sex, pregnancy, specimen type, assay, and laboratory; a result outside a range is not automatically a diagnosis, and a result within a range does not always exclude disease. Discuss unexpected results with the ordering clinician, who can decide whether repeat or confirmatory testing, treatment, or referral is appropriate. Do not change medicines or delay urgent care based on an isolated result.


\section*{Regulatory Status and Authorization Date}
Regulatory status applies to the specific assay, instrument, manufacturer, and intended use rather than to the test name alone. No single approval date can be assigned to this test category. For a commercial assay, verify the current product in the relevant regulator's database: the U.S. Food and Drug Administration (FDA) clearance, approval, or authorization; India's Central Drugs Standard Control Organisation (CDSCO) licence or permission; the European Union In Vitro Diagnostic Regulation (EU IVDR) CE certificate issued through the applicable conformity-assessment route; or the United Kingdom Medicines and Healthcare products Regulatory Agency (MHRA) registration and UKCA/CE status. Laboratory-developed tests may not have an individual FDA, CDSCO, EU IVDR, or MHRA approval date.
\clearpage
